%======================================================================
%	Vorlage
%======================================================================
%	$Id$
%	Matthias Kupfer
%======================================================================
%	Documentclass
%======================================================================
\documentclass[
%	draft,			% Entwurfsmodus: Bilder als Rahmen,
				% Überlängen werden deutlich markiert
%	10pt,			% 
	11pt,			% KOMA default
%	12pt,			% 
	a4paper,		% DIN A4
%	twoside,		% Zweiseitig
%	german,			%
    english
%----------------------------------------------------------------------
% Die folgenden Befehle stammen aus dem KOMA Paket
	headsepline,		% Linie unter der Kopfzeile	
%	headnosepline,		%
%	foodsepline,		% Linie über Fussnote	
	footsepline=false,	%
	automark,		% Kolumnentitel lebendig
%	bigheadings,		% default: Überschriften gross setzen
	headings=normal,	% Überschriften normal setzen
%	smallheadings,		% Überschriften eher klein setzen
	pointlessnumbers,	% Keinen Punkt hinter die letzte Zahl 
				% eines Kapitels (auch bei Anhang)
%	chapterprefix,		% Kapiteluberschriften "Kapitel"
%	appendixprefix,		% Anhang
%	openright,		% Kapitel auf der rechten (ungeraden) Seite 
				% anfangen lassen
	openany,		% 
%	cleardoublestandard,	%
	cleardoublepage=plain,	%
%	cleardoubleempty,	%
	abstracton,		%
	index=totoc,		% Index soll im Inhaltsverzeichnis auftauchen
	listof=totoc,		%
	bibliography=totoc,	%
%	parskip, 		% parskip-, parskip*, parskip+
% 	halfparskip, 		% halfparskip-, halfparskip* und halfparskip+
%	DIVclassic,
% 	BCOR=8mm,		%%?? Bindungskorrektur: 
				% BCOR<Breite des Bindungsverlustes> 
]{scrreprt} 

%======================================================================
%	Bearbeitung sowohl mit LaTeX als auch mit pdfLaTeX ermoeglichen
%======================================================================
%\newif\ifpdf 
	%\ifx\pdfoutput\undefined 
	%\pdffalse 			% we are not running pdflatex 
%\else 
	%\pdfoutput=1 			% we are running pdflatex 
	%\pdftrue 
%\fi


%======================================================================
%	Verwendete Pakete
%======================================================================

%\usepackage{babel}		% Sprachen

%% Latex mit deutschen Umlauten:
%% http://www.cs.albany.edu/~herrmann/latex_umlaute/
\usepackage[T1]{fontenc}	% EC-Schriften verwenden (vs. DC) da 8-Bit
				% EC-Schriften als T1-kodierten CM-Schriften
				% European/Ext.-Computer-Modern-(EC)-Schriften
				% Umlaute, Anführungszeichen ...
				% => Umlauten koennen richtig getrennt werden
				% FAQ 5.3.2

%\usepackage{ucs}		% Eingabe von ü,ö,ä,ß mit UTF-8 erlaubt
%\usepackage[utf8]{inputenc}	% (auch in Include-Dateien)
				% ! Zeichensatz der Dateien als UTF-8 !
								
\usepackage{ae,aecompl}		% virtuelle-CM-Fonts
				% da EC nicht als PostScript-(Type-1) verfuegbar
				% => keine echten Umlaute im PDF-Dokumen 
				%(Problem bei Suche)
				% By loading the ae package (\usepackage{ae}), 
				% you loose some characters as mentioned in 
				% README. 
				% The package aecompl by Denis Roegel restores
				% these characters which are taken from the ec 
				% fonts. If you use pdftex, you will get these 
				% characters as bitmaps, but this might be 
				% better than not having them at all.
\usepackage{bbold}
\usepackage{amsmath}
\usepackage[hang]{caption}
%\usepackage{times, mathptm}	% TimesNewRoman Schrift (Acrobat Reader Fonts), 
				% dazu braucht man auch den entsprechenden 
				% Zeichensatz für den Math-Mode
%\usepackage{pslatex}		% ? mathematische Formeln mit Standard 
				% Postscript Fonts gesetzt
			% Paket     Roman         Serifenlos  Typewriter
%\usepackage{times}	% -----------------------------------------------
			% times     Times         Helvetica   Courier
			% palatino  Palatino      Helvetica   Courier
			% newcent   NewCenturySch AvantGarde  Courier
			% bookman   Bookman       AvantGarde  Courier
			% Diese Schriften sind die Standard-PostScript-Schriften
			% und in jedem Drucker verfügbar

\usepackage{
%	german,			% Deutsche Trennungen (ALTE Rechtschreibung), 
				% Anführungsstriche und mehr 
%	ngerman,		% Deutsche Trennungen (NEUE Rechtschreibung), 
				% Anführungsstriche und mehr
				%
%	acronym,		% Verwaltung von Abkuerzungen
%	bibgerm,		% Deutsche Bibliographie 
	calc,			% Erweiterung der arithmetischen Funktionen in 
				% LaTeX
				% wird verwendet um Titelseite zu zentrieren
	color,			% im Laufenden Text einfach mit \color{Farbe) zwischen den 
				% Farben umschalten, wobei Farbe einfach 
				% durch z.B. red, blue, black etc. ersetzt wird
				% \textcolor{farbe){Text)
	datetime,   % definiere Datumsstile
	enumitem,   % mehr Optionen für enumerate environments 
%	epigraph,		% Zitat am Kapitelanfang
%	fancyhdr,		% Kopf- und Fußzeilen von Dokumenten frei 
				% gestalten
	fancybox,		% shadowbox, doublebox, ovalbox, Ovalbox 
	fancyvrb,		% verbatim Erweiterung:
	float,			% Positionierung von Gleitobjekten genau an der Stelle, wo man
				% 'figure'- oder 'table'-Umgebung die 
				% Positionierung [H] gesetzt werden
	%glossary,		% Glossar und Abkürzungsverzeichnis
	mdwlist,		% compact list: itemize* ..
	scrdate,		% \todaysname 
	scrtime,		% \thistime
	scrlayer-scrpage,		% Kopf- und Fußzeilen flexibel gestalten
				%
%	moreverb,		% verbatim-ähnlich: boxedverbatim, listing
%	verbatim,		% Darstellung von "Text, wie er eingegeben wird"
				%
%	lscape,			% Erstellt eine um 90% gedrehte *neue* Seite
%	textcomp,		% Sonderzeichen
%	booktabs,		% Tabellenlinien
%	longtable,		% Tabellen > 1 Seite
%	supertabular,		% Tabellen > 1 Seite
	tabularx,		% Blocksatzspalten
%	ltxtable,		% tabularx + longtable
%	multicol,		% mehrspaltige Zeilen
%	varioref,		% einheitliche Verweise
%	endnotes,		% Fussnoten -> Endnoten
%	rotating,		% sidewaystable und sidewaysfigure
%	natbib,			% Bibliographie ohne Klammer etc.
%	marvosym,		% Euro etc.
}


\usepackage[acronym,toc]{glossaries}
\makeglossaries


\usepackage{listings}

%\usepackage{pstricks}
%\usepackage{listings}
%\usepackage{wasysym}

%\usepackage[german, first, bottomafter]{draftcopy}

%\usepackage{setspace}	% Durchschuß, Zeilenabstand
%\doublespace		% doppelzeilig oder
%\onehalfspacing	% anderthalbzeilig
\renewcommand{\baselinestretch}{1.1}\normalsize

% Für schöne Darstellung von Algorithmen
%\usepackage[german, algoruled, algochapter]{algorithm2e}
%\usepackage{algorithmic}
%\usepackage[chapter]{algorithm}
%\floatname{algorithm}{Algorithmus}

%======================================================================
%	Pictures, Links
%======================================================================

\usepackage{graphicx}
\graphicspath{Pictures/}	% Angabe der Pfade, wo die Grafiken liegen; 
				% mehrere Pfade sind möglich
\usepackage{subfig}
\usepackage{thumbpdf}
\usepackage[breaklinks]{hyperref}
\def\UrlBreaks{\do\/\do-}
%\usepackage{breakurl}
\usepackage{minted}



%======================================================================
%	user.sty
%======================================================================

\usepackage{user}	% Makros in user.sty 
			% \epsinc{bild.eps}{scale=1}{Bildunterschrift}
			% \missing{Beweis fehlt noch}
			% \comment{Ein Kommentar}

%======================================================================
%	Settings
%======================================================================

%\typearea{11}		% Satzspiegel neu konstruieren (KOMA)
			%      10pt 11pt 12pt 
			% DIV : 8   10   12  

\pagestyle{scrheadings}	% Standart  Kopf- und Fußzeile
\setkomafont{pageheadfoot}{\small\scshape}

% ---------------------------------------------------------------------
%\setcounter{secnumdepth}{2}
%\setcounter{chapter}{-1}
\setcounter{tocdepth}{3}

% ---------------------------------------------------------------------
%\sloppy		% weniger Worttrennungen, größere Wortabstände
\fussy			% viele Worttrennungen, "schönere" Wortabstände

% ---------------------------------------------------------------------
%\flushbottom			% Ausrichtung der Seitenenden jeweils auf 
				% gleicher Höhe

% ---------------------------------------------------------------------
%\sloppypar		% Das hier relaxt die Einstellungen zum Wortabstand 
			% extrem. Damit ragen keine Worte über den rechten 
			% Zeilenabstand hinaus. Dafür muß stärker auf
			% Wortabstand geachtet werden, der kann dann ziemlich 
			% groß werden. Man erhält aber keine Meldung mehr 
			% über underfull boxes.

%Hiermit kann man das gleiche mit weniger Holzhammer erreichen:
%\setlength{\tolerance}{2000}           % Strafpunkt für Zeilenumbruch
%\setlength{\emergencystretch}{3pt}     % Soweit dürfen einzelne Worte mehr 
					% auseinandergezogen werden
%\setlength{\hfuzz}{1pt}                % Macht den rechten Rand um bis zu 1pt 
					% flatterig.

% ---------------------------------------------------------------------
%Vermeiden einzelner Zeilen am Ende einer Seite oder oben auf einer neuen Seite
%\clubpenalty10000
%\widowpenalty10000

% ----------------------------------------------------------------------
% neue Umgebungen für verwendete Sätze und Beispiele
\newtheorem{bsp}{Beispiel}[chapter]
\newtheorem{satz}{Satz}[chapter]

% Eigene Anpassungen
% Silbentrennung
\tolerance 1414
\hbadness 1414                                                                  \emergencystretch 1.5em
\hfuzz 0.3pt
\widowpenalty=10000
\hyphenpenalty=200
\clubpenalty = 10000
\widowpenalty = 10000
\displaywidowpenalty = 10000
\vfuzz \hfuzz
\raggedbottom


%======================================================================
%	includeonly
%======================================================================

\includeonly{		% Shows which file shall be included by the include command
		%
  metadata		% set in variables
  ,title		% title page, Summary and table of contents
% All files for the include command have to be listed here!
  ,Text/Introduction	%
  ,Text/Fundamentals
  ,Text/Implementation
  ,Text/Results
  ,Text/Evaluation
  ,Text/Summary_Outlook
}

% Recommended structure of your scientific work
%
% Title Page
% Bibliographic Description (Backside of Title Page)
% (Task)
% Acknowledgement
% Abstract
% Table of Contents
% List of Tables
% List of Figures
% Introduction
% Content
% Summary/Outlook
% (List of used terms)
% Glossary (List of Abbreviations)
% (Index)
% Bibliography/References
% (Theses)
% (Declaration of Authorship)
% Attachments


%======================================================================
% * *   T E X T    S T A R T S    H E R E   * * * * * * * * * * * * * *
%======================================================================

\begin{document}
%======================================================================
%	Metadata
%======================================================================
%	$Id$
%	Matthias Kupfer
%======================================================================

\newcommand{\dcsubject}{Master Thesis}
% z.B. (Diplom/Studien/Haus)arbeit, Praktikumsbericht, Studie, Beleg, 
% (Pro/Haupt/Ober)seminar, Seminar usw.
\newcommand{\dctitle}{System Design, Implementation and Evaluation of a User-centered Application for Indoor Sound Classification}
\newcommand{\dcsubtitle}{~} % Subtitle, if necessary

\newcommand{\dcauthorlastname}{Sun}
\newcommand{\dcauthorfirstname}{Maoyin}
\newcommand{\dcauthoremail}{maoyin.sun@s2016.tu-chemnitz.de}
\newdateformat{mydate}{{\THEDAY} \monthname[\THEMONTH] \THEYEAR}
\newcommand{\dcdate}{\mydate\today}
%\newcommand{\dcdate}{\today}

\newcommand{\dcplace}{Chemnitz} % Ort, kann an der TU meist so bleiben
\newcommand{\dcuni}{Technische Universität \dcplace}
\newcommand{\dcdepart}{Faculty of Computer Science} % Fakultätsangabe
\newcommand{\dcprof}{Junior Professorship of Media Computing} % Angabe der Professur

\newcommand{\dcpruefer}{Jun.-Prof. Dr. rer. nat. Danny Kowerko}% Prüfer der Arbeit
\newcommand{\dcadvisor}{M. Sc. Arunodhayan Sampath-Kumar}% Betreuer der Arbeit

\newcommand{\dckeywords}{System Design, User-centered Application, Indoor Sound Classification, Web-based Demonstration Platform, Asynchronous Messaging, Acoustic Event Classification, Decoupled Architecture, MQTT, D3.js, Real-time Streaming}

%%======================================================================
% Settings of Hyperref-Package
\hypersetup{%    
    pdftitle	= {\dctitle}, %
    pdfsubject	= {\dcsubject}, %
    pdfauthor	= {\dcauthorfirstname~\dcauthorlastname, \dcauthoremail}, %
    %	pdfkeywords	= {\dckeywords}, %
    pdfcreator	= {pdfTeX with hyperref and Thumbpdf}, %
}
%%======================================================================
 % Variables, \hypersetup, etc.

%======================================================================
%	Title Page (Abstract and Table of Contents)
%======================================================================

%======================================================================
%	Title Page
%======================================================================
%	$Date:$
%	$Revision:$
%	Matthias Kupfer
%======================================================================

%%======================================================================
%% Extra Title
%%======================================================================
%\extratitle{
%	\usekomafont{disposition}\mdseries 
%	\begin{center}
%		\Huge \dcsubject\\[1.5ex]
%		\hrule
%		\vspace*{\fill}
%		\includegraphics{TUC_deutsch_einzeile_CMYK}
%	\end{center}
%}

%%======================================================================
%% Title Head
%%======================================================================
\titlehead{
  \vspace*{-1.5cm}
  % font families like all headings, but not bold
  \usekomafont{disposition}\mdseries
  \begin{center}
    \raisebox{-1ex}{\includegraphics[scale=1.4]{TUC_deutsch_einzeile_CMYK}}\\
    \hrulefill \\[1em]
    {\Large\dcdepart}\\[0.5em]
    \dcprof
  \end{center}
  \vspace*{1.5cm}
}

%%======================================================================
%% Subject
%%======================================================================


\subject{\normalfont\bfseries\Huge\dcsubject}

%%======================================================================
%% Title
%%======================================================================
\title{\normalfont\sffamily\Large
  \dctitle
  \\
  \dcsubtitle
}

%%======================================================================
%% Author of Document
%%======================================================================
\author{\dcauthorfirstname~\dcauthorlastname}

%%======================================================================
%% Place, Date
%%======================================================================
\date{\dcplace, \dcdate}

%%======================================================================
%% Publishers
%%======================================================================
\publishers{
  {\parbox{\textwidth-8em}{
        \begin{tabbing}
          {\normalfont\bfseries Supervisor:}\quad\=\kill
          {\normalfont\bfseries Examiner:}	\>\dcpruefer\\
          {\normalfont\bfseries Supervisor:}	\>\dcadvisor
        \end{tabbing}
      }}
}

%%======================================================================
%% Bibliographic Statement
%%======================================================================
\lowertitleback{
  \textbf{\dcauthorlastname, \dcauthorfirstname}\\
  \dctitle\\
  \dcsubject,~\dcdepart\\
  \dcuni,~\ifcase\month\or
    January\or February\or March\or April\or May\or June\or
    July\or August\or September\or October\or November\or December\fi
  ~\number\year
}

%%======================================================================
%% maketitle
%%======================================================================

\maketitle



%%======================================================================
%% Acknowledgement
%%======================================================================
\thispagestyle{empty}
%\null\vfil
\begin{center}
  \usekomafont{disposition}
  \textbf{{Acknowledgement}}
  %\vspace{-.5em}\vspace{\parsep}
\end{center}

I would like to express my sincere gratitude to a number of people whose support and guidance have been invaluable throughout the journey of this thesis.

First and foremost, my deepest appreciation goes to my thesis advisors, Prof. Danny Kowerko and Mr. Arunodhayan Sampath-Kumar, whose expertise, understanding, and patience, added considerably to my graduate experience. Your guidance was invaluable at every stage of this research, and I am immensely grateful for your input and encouragement.

I also wish to express my appreciation to the faculty and staff in the Junior Professorship of Media Computing at TU Chemnitz. Your dedication to fostering a rigorous academic environment has profoundly enriched my educational experience.

I would also like to acknowledge the technical support and resources provided by our university, which were crucial to the successful completion of this project.

To my family, thank you for your unwavering support and encouragement throughout my years of study and through the process of researching and writing this thesis. This accomplishment would not have been possible without you.

Lastly, I cannot forget to express my gratitude to all those who were involved in the practical aspects of this project. Your contributions have been immensely valuable.

Thank you all for your support and encouragement.

\let\cleardoublepage\cleardoubleemptypage
%\cleardoublepage

%%======================================================================
%%       Abstract
%%======================================================================
\begin{abstract}
  The increasing demand for robust and efficient acoustic event classification models necessitates innovative approaches to their demonstration and deployment. This thesis presents a novel web-based demonstration platform designed to showcase the capabilities of an acoustic event classification model proposed by \cite{sampath2019cnn}. Emphasizing user-friendliness, configurability, scalability, and robustness, the implementation leverages asynchronous messaging and a decoupled architecture, utilizing MQTT as the messaging protocol for the backend and \textsc{d3.js} for visualization in the frontend.

  The user-friendly interface allows users to interact intuitively with the acoustic event classification model, facilitating a seamless experience. The system's configurability empowers users to adapt the demonstration to various scenarios and applications. Scalability is achieved through a decoupled architecture that ensures efficient resource utilization and accommodates varying workloads. The thesis explores the integration of asynchronous messaging to enhance the responsiveness and real-time capabilities of the demonstration. By decoupling components, the architecture not only improves scalability but also contributes to the system's robustness, ensuring reliable performance under diverse conditions. Throughout the thesis, we detail the implementation and deployment of this demonstration platform, emphasizing the technical considerations of the asynchronous messaging system and the advantages of a decoupled architecture. Evaluation results demonstrate the platform's efficiency, highlighting its potential for broader applications in acoustic event classification.

  This research contributes to the intersection of acoustic event classification, web-based demonstrations, and system architecture, offering a valuable resource for researchers, practitioners, and developers in the field. The demonstrated platform serves as a benchmark for user-friendly and scalable acoustic event classification models, paving the way for advancements in real-world applications.

\end{abstract}

%%======================================================================
%%      Table of Contents
%%======================================================================
%\cleardoubleemptypage
\pdfbookmark{Contents}{Contents}
\pagenumbering{roman}
\tableofcontents

%%======================================================================
%%      List of Figures
%%======================================================================
%\cleardoublepage
%\markboth{List of Figures}{List of Figures}
\listoffigures

%%======================================================================
%%      List of Tables
%%======================================================================
%\cleardoublepage
%\markboth{List of Tables}{List of Tables}
\listoftables

%%======================================================================
%%      List of Algorithms
%%======================================================================
%\renewcommand{\listalgorithmname}{List of Algorithms}
%\cleardoublepage
%\addcontentsline{toc}{chapter}{List of Algorithms}
%\listofalgorithms

%%======================================================================
%%      Glossary
%%======================================================================
%\cleardoublepage               

\newacronym{acronym1}{ACR1}{Acronym 1}
\newacronym{acronym2}{ACR2}{Acronym 2}
\newacronym{acronym3}{ACR3}{Acronym 3}

\printglossary
\printglossary[type=\acronymtype,title=List of Acronyms,toctitle=List of Acronyms]

%%======================================================================
%%      End
%%======================================================================
%\cleardoublepage
%\pagenumbering{arabic}

 % Titelpage, Abstract 
		% Table of Contents
		% Glossary
		% -> then page numbering starting at 1

%======================================================================
%	Content per include
%======================================================================

\chapter{Introduction}
\setcounter{page}{1}
\pagenumbering{arabic}

\section{Context}
Lorem ipsum dolor sit amet \acrfull{acronym1}, consetetur sadipscing elitr, sed diam nonumy eirmod tempor invidunt ut labore et dolore magna aliquyam erat, sed diam voluptua. At vero eos et accusam et justo duo dolores et ea rebum. Stet clita kasd gubergren, no sea takimata sanctus est Lorem ipsum dolor sit amet. \textit{Lorem ipsum} dolor sit amet, consetetur sadipscing elitr, sed diam nonumy eirmod tempor invidunt ut labore et dolore magna aliquyam erat, sed diam voluptua. At vero eos et accusam et justo duo dolores et ea rebum. Stet clita kasd gubergren, no sea takimata sanctus est Lorem ipsum dolor sit amet. Lorem ipsum dolor sit amet, consetetur sadipscing elitr, sed diam nonumy eirmod tempor invidunt ut labore et dolore magna aliquyam erat, sed diam voluptua. At vero eos et accusam et justo duo dolores et ea rebum. Stet clita kasd gubergren, no sea takimata sanctus est Lorem ipsum dolor sit amet \cite{KHF17_Acoustic_Event_Classification_Using_Convolutional_Neural_Networks}. 

\section{Motivation}
Duis autem vel eum iriure dolor in hendrerit in vulputate velit esse molestie consequat, vel illum dolore eu \acrfull{acronym2} feugiat nulla facilisis at vero eros et accumsan et iusto odio dignissim qui blandit praesent luptatum zzril delenit augue duis dolore te feugait nulla facilisi. Lorem ipsum dolor sit amet, consectetuer adipiscing elit, sed diam nonummy nibh euismod tincidunt ut laoreet dolore magna aliquam erat volutpat \cite{Pic15_Environmental_Sound_Classification_with_Convolutional_Neural_Networks}. 

\section{Aims of the work and methodical approach}
List  central theses of your thesis title derived from context and motivation:

Ut wisi enim ad minim veniam, quis nostrud exerci tation ullamcorper suscipit lobortis nisl ut aliquip ex ea commodo consequat. Duis autem vel eum iriure dolor in hendrerit in vulputate velit esse molestie consequat, vel illum dolore eu feugiat nulla facilisis at vero eros et \acrfull{acronym3} accumsan et iusto odio dignissim qui blandit praesent luptatum zzril delenit augue duis dolore te feugait nulla facilisi \cite{SAP17_Unsupervised_Filterbank_Learning_Using_Convolutional_Restricted_Boltzmann}.

\section{Outline of the Thesis}





\chapter{Fundamentals and state of the art}
\section{Section\_1}

Lorem ipsum dolor sit amet \acrfull{acronym1}, consetetur sadipscing elitr, sed diam nonumy eirmod tempor invidunt ut labore et dolore magna aliquyam erat, sed diam voluptua. At vero eos et accusam et justo duo dolores et ea rebum. Stet clita kasd gubergren, no sea takimata sanctus est Lorem ipsum dolor sit amet. Lorem ipsum dolor sit amet, consetetur sadipscing elitr, sed diam nonumy eirmod tempor invidunt ut labore et dolore magna aliquyam erat, sed diam voluptua. At vero eos et accusam et justo duo dolores et ea rebum. Stet clita kasd gubergren, no sea takimata sanctus est Lorem ipsum dolor sit amet. Lorem ipsum dolor sit amet, consetetur sadipscing elitr, sed diam nonumy eirmod tempor invidunt ut labore et dolore magna aliquyam erat, sed diam voluptua. At vero eos et accusam et justo duo dolores et ea rebum. Stet clita kasd gubergren, no sea takimata sanctus est Lorem ipsum dolor sit amet \cite{fedden_comparative_nodate}. 

Duis autem vel eum iriure dolor in hendrerit in vulputate velit esse molestie consequat, vel illum dolore eu \acrfull{acronym2} feugiat nulla facilisis at vero eros et accumsan et iusto odio dignissim qui blandit praesent luptatum zzril delenit augue duis dolore te feugait nulla facilisi. Lorem ipsum dolor sit amet, consectetuer adipiscing elit, sed diam nonummy nibh euismod tincidunt ut laoreet dolore magna aliquam erat volutpat. 

Ut wisi enim ad minim veniam, quis nostrud exerci tation ullamcorper suscipit lobortis nisl ut aliquip ex ea commodo consequat. Duis autem vel eum iriure dolor in hendrerit in vulputate velit esse molestie consequat, vel illum dolore eu feugiat nulla facilisis at vero eros et \acrfull{acronym3} accumsan et iusto odio dignissim qui blandit praesent luptatum zzril delenit augue duis dolore te feugait nulla facilisi.

\section{Section\_2}

Lorem ipsum dolor sit amet \acrfull{acronym1}, consetetur sadipscing elitr, sed diam nonumy eirmod tempor invidunt ut labore et dolore magna aliquyam erat, sed diam voluptua. At vero eos et accusam et justo duo dolores et ea rebum. Stet clita kasd gubergren, no sea takimata sanctus est Lorem ipsum dolor sit amet. Lorem ipsum dolor sit amet, consetetur sadipscing elitr, sed diam nonumy eirmod tempor invidunt ut labore et dolore magna aliquyam erat, sed diam voluptua. At vero eos et accusam et justo duo dolores et ea rebum. Stet clita kasd gubergren, no sea takimata sanctus est Lorem ipsum dolor sit amet. Lorem ipsum dolor sit amet, consetetur sadipscing elitr, sed diam nonumy eirmod tempor invidunt ut labore et dolore magna aliquyam erat, sed diam voluptua. At vero eos et accusam et justo duo dolores et ea rebum. Stet clita kasd gubergren, no sea takimata sanctus est Lorem ipsum dolor sit amet. 

Duis autem vel eum iriure dolor in hendrerit in vulputate velit esse molestie consequat, vel illum dolore eu \acrfull{acronym2} feugiat nulla facilisis at vero eros et accumsan et iusto odio dignissim qui blandit praesent luptatum zzril delenit augue duis dolore te feugait nulla facilisi. Lorem ipsum dolor sit amet, consectetuer adipiscing elit, sed diam nonummy nibh euismod tincidunt ut laoreet dolore magna aliquam erat volutpat. 

Ut wisi enim ad minim veniam, quis nostrud exerci tation ullamcorper suscipit lobortis nisl ut aliquip ex ea commodo consequat. Duis autem vel eum iriure dolor in hendrerit in vulputate velit esse molestie consequat, vel illum dolore eu feugiat nulla facilisis at vero eros et \acrfull{acronym3} accumsan et iusto odio dignissim qui blandit praesent luptatum zzril delenit augue duis dolore te feugait nulla facilisi.

\section{Section\_3}

Lorem ipsum dolor sit amet \acrfull{acronym1}, consetetur sadipscing elitr, sed diam nonumy eirmod tempor invidunt ut labore et dolore magna aliquyam erat, sed diam voluptua. At vero eos et accusam et justo duo dolores et ea rebum. Stet clita kasd gubergren, no sea takimata sanctus est Lorem ipsum dolor sit amet. Lorem ipsum dolor sit amet, consetetur sadipscing elitr, sed diam nonumy eirmod tempor invidunt ut labore et dolore magna aliquyam erat, sed diam voluptua. At vero eos et accusam et justo duo dolores et ea rebum. Stet clita kasd gubergren, no sea takimata sanctus est Lorem ipsum dolor sit amet. Lorem ipsum dolor sit amet, consetetur sadipscing elitr, sed diam nonumy eirmod tempor invidunt ut labore et dolore magna aliquyam erat, sed diam voluptua. At vero eos et accusam et justo duo dolores et ea rebum. Stet clita kasd gubergren, no sea takimata sanctus est Lorem ipsum dolor sit amet. Duis autem vel eum iriure dolor in hendrerit in vulputate velit esse molestie consequat, vel illum dolore eu \acrfull{acronym2} feugiat nulla facilisis at vero eros et accumsan et iusto odio dignissim qui blandit praesent luptatum zzril delenit augue duis dolore te feugait nulla facilisi:

\begin{itemize}
    \item \underline{Item\_1:} 
    \begin{itemize}
        \item Subitem\_1, Subitem\_2, Subitem\_3.
    \end{itemize}

    \item \underline{Item\_1:} 
    \begin{itemize}
        \item Subitem\_1, Subitem\_2, Subitem\_3.
    \end{itemize}

    \item \underline{Item\_3:}  
    \begin{itemize}
        \item Subitem\_1, Subitem\_2, Subitem\_3.
    \end{itemize}
\end{itemize}

\chapter{Implementation and methods}

Lorem ipsum dolor sit amet \acrfull{acronym1}, consetetur sadipscing elitr, sed diam nonumy eirmod tempor invidunt ut labore et dolore magna aliquyam erat, sed diam voluptua. At vero eos et accusam et justo duo dolores et ea rebum. Stet clita kasd gubergren, no sea takimata sanctus est Lorem ipsum dolor sit amet. Lorem ipsum dolor sit amet, consetetur sadipscing elitr, sed diam nonumy eirmod tempor invidunt ut labore et dolore magna aliquyam erat, sed diam voluptua. At vero eos et accusam et justo duo dolores et ea rebum. Stet clita kasd gubergren, no sea takimata sanctus est Lorem ipsum dolor sit amet. Lorem ipsum dolor sit amet, consetetur sadipscing elitr, sed diam nonumy eirmod tempor invidunt ut labore et dolore magna aliquyam erat, sed diam voluptua. At vero eos et accusam et justo duo dolores et ea rebum. Stet clita kasd gubergren, no sea takimata sanctus est Lorem ipsum dolor sit amet.

\section{Section\_1}

Lorem ipsum dolor sit amet \acrfull{acronym1}, consetetur sadipscing elitr, sed diam nonumy eirmod tempor invidunt ut labore et dolore magna aliquyam erat, sed diam voluptua. At vero eos et accusam et justo duo dolores et ea rebum. Stet clita kasd gubergren, no sea takimata sanctus est Lorem ipsum dolor sit amet. Lorem ipsum dolor sit amet, consetetur sadipscing elitr, sed diam nonumy eirmod tempor invidunt ut labore et dolore magna aliquyam erat, sed diam voluptua. At vero eos et accusam et justo duo dolores et ea rebum. Stet clita kasd gubergren, no sea takimata sanctus est Lorem ipsum dolor sit amet. Lorem ipsum dolor sit amet, consetetur sadipscing elitr, sed diam nonumy eirmod tempor invidunt ut labore et dolore magna aliquyam erat, sed diam voluptua. At vero eos et accusam et justo duo dolores et ea rebum. Stet clita kasd gubergren, no sea takimata sanctus est Lorem ipsum dolor sit amet. 

Duis autem vel eum iriure dolor in hendrerit in vulputate velit esse molestie consequat, vel illum dolore eu \acrfull{acronym2} feugiat nulla facilisis at vero eros et accumsan et iusto odio dignissim qui blandit praesent luptatum zzril delenit augue duis dolore te feugait nulla facilisi. Lorem ipsum dolor sit amet, consectetuer adipiscing elit, sed diam nonummy nibh euismod tincidunt ut laoreet dolore magna aliquam erat volutpat. 
\ref{fig:Specs}

Ut wisi enim ad minim veniam, quis nostrud exerci tation ullamcorper suscipit lobortis nisl ut aliquip ex ea commodo consequat. Duis autem vel eum iriure dolor in hendrerit in vulputate velit esse molestie consequat, vel illum dolore eu feugiat nulla facilisis at vero eros et \acrfull{acronym3} accumsan et iusto odio dignissim qui blandit praesent luptatum zzril delenit augue duis dolore te feugait nulla facilisi.

\begin{figure}[H]
\frame{\includegraphics[width=0.475\textwidth]{Pictures/1-71030-A-6_1_hen.png}}
\hspace{0.05\textwidth}
\frame{\includegraphics[width=0.475\textwidth]{Pictures/2-117617-A-48_1_fireworks.png}}\\\\
\frame{\includegraphics[width=0.475\textwidth]{Pictures/3-155578-A-14_1_chirping_birds.png}}
\hspace{0.05\textwidth}
\frame{\includegraphics[width=0.475\textwidth]{Pictures/3-159347-A-36_0_vaccumcleaner.png}}\\\\
\frame{\includegraphics[width=0.475\textwidth]{Pictures/4-197454-B-28_1_snoring.png}}
\hspace{0.05\textwidth}
\frame{\includegraphics[width=0.475\textwidth]{Pictures/1-76831-C-42_2_siren.png}}
\caption{Ut wisi enim ad minim veniam, quis nostrud exerci tation ullamcorper suscipit lobortis nisl ut aliquip ex ea commodo consequat. Duis autem vel eum iriure dolor in hendrerit in vulputate velit esse molestie consequat, vel illum dolore eu feugiat nulla facilisis at vero eros et \acrfull{acronym3} accumsan et iusto odio dignissim qui blandit praesent luptatum zzril delenit augue duis dolore te feugait nulla facilisi.}
\label{fig:Specs}
\end{figure}

\section{Section\_2}

Ut wisi enim ad minim veniam, quis nostrud exerci tation ullamcorper suscipit lobortis nisl ut aliquip ex Appendix \ref{chap:appendix} ea commodo consequat.

\subsection{Subsection\_1}

Lorem ipsum dolor sit amet \acrfull{acronym1}, consetetur sadipscing elitr, sed diam nonumy eirmod tempor invidunt ut labore et dolore magna aliquyam erat, sed diam voluptua. At vero eos et accusam et justo duo dolores et ea rebum. Stet clita kasd gubergren, no sea takimata sanctus est Lorem ipsum dolor sit amet. Lorem ipsum dolor sit amet, consetetur sadipscing elitr, sed diam nonumy eirmod tempor invidunt ut labore et dolore magna aliquyam erat, sed diam voluptua. At vero eos et accusam et justo duo dolores et ea rebum. Stet clita kasd gubergren, no sea takimata sanctus est Lorem ipsum dolor sit amet. Lorem ipsum dolor sit amet, consetetur sadipscing elitr, sed diam nonumy eirmod tempor invidunt ut labore et dolore magna aliquyam erat, sed diam voluptua. At vero eos et accusam et justo duo dolores et ea rebum. Stet clita kasd gubergren, no sea takimata sanctus est Lorem ipsum dolor sit amet.

\subsection{Subsection\_2}

Lorem ipsum dolor sit amet \acrfull{acronym1}, consetetur sadipscing elitr, sed diam nonumy eirmod tempor invidunt ut labore et dolore magna aliquyam erat, sed diam voluptua. At vero eos et accusam et justo duo dolores et ea rebum. Stet clita kasd gubergren, no sea takimata sanctus est Lorem ipsum dolor sit amet. Lorem ipsum dolor sit amet, consetetur sadipscing elitr, sed diam nonumy eirmod tempor invidunt ut labore et dolore magna aliquyam erat, sed diam voluptua. At vero eos et accusam et justo duo dolores et ea rebum. Stet clita kasd gubergren, no sea takimata sanctus est Lorem ipsum dolor sit amet. Lorem ipsum dolor sit amet, consetetur sadipscing elitr, sed diam nonumy eirmod tempor invidunt ut labore et dolore magna aliquyam erat, sed diam voluptua. At vero eos et accusam et justo duo dolores et ea rebum. Stet clita kasd gubergren, no sea takimata sanctus est Lorem ipsum dolor sit amet.


\chapter{Results and discussion}

Lorem ipsum dolor sit amet \acrfull{acronym1}, consetetur sadipscing elitr, sed diam nonumy eirmod tempor invidunt ut labore et dolore magna aliquyam erat, sed diam voluptua. At vero eos et accusam et justo duo dolores et ea rebum. Stet clita kasd gubergren, no sea takimata sanctus est Lorem ipsum dolor sit amet. Lorem ipsum dolor sit amet, consetetur sadipscing elitr, sed diam nonumy eirmod tempor invidunt ut labore et dolore magna aliquyam erat, sed diam voluptua. At vero eos et accusam et justo duo dolores et ea rebum. Stet clita kasd gubergren, no sea takimata sanctus est Lorem ipsum dolor sit amet. Lorem ipsum dolor sit amet, consetetur sadipscing elitr, sed diam nonumy eirmod tempor invidunt ut labore et dolore magna aliquyam erat, sed diam voluptua. At vero eos et accusam et justo duo dolores et ea rebum. Stet clita kasd gubergren, no sea takimata sanctus est Lorem ipsum dolor sit amet. 

Duis autem vel eum iriure dolor in hendrerit in vulputate velit esse molestie consequat, vel illum dolore eu \acrfull{acronym2} feugiat nulla facilisis at vero eros et accumsan et iusto odio dignissim qui blandit praesent luptatum zzril delenit augue duis dolore te feugait nulla facilisi. Lorem ipsum dolor sit amet, consectetuer adipiscing elit, sed diam nonummy nibh euismod tincidunt ut laoreet dolore magna aliquam erat volutpat. 

Ut wisi enim ad minim veniam, quis nostrud exerci tation ullamcorper suscipit lobortis nisl ut aliquip ex ea commodo consequat. Duis autem vel eum iriure dolor in hendrerit in vulputate velit esse molestie consequat, vel illum dolore eu feugiat nulla facilisis at vero eros et \acrfull{acronym3} accumsan et iusto odio dignissim qui blandit praesent luptatum zzril delenit augue duis dolore te feugait nulla facilisi.

\section{Section\_1}

Lorem ipsum dolor sit amet \acrfull{acronym1}, consetetur sadipscing elitr, sed diam nonumy eirmod tempor invidunt ut labore et dolore magna aliquyam erat, sed diam voluptua. At vero eos et accusam et justo duo dolores et ea rebum. Stet clita kasd gubergren, no sea takimata sanctus est Lorem ipsum dolor sit amet. Lorem ipsum dolor sit amet, consetetur sadipscing elitr, sed diam nonumy eirmod tempor invidunt ut labore et dolore magna aliquyam erat, sed diam voluptua. At vero eos et accusam et justo duo dolores et ea rebum. Stet clita kasd gubergren, no sea takimata sanctus est Lorem ipsum dolor sit amet. Lorem ipsum dolor sit amet, consetetur sadipscing elitr, sed diam nonumy eirmod tempor invidunt ut labore et dolore magna aliquyam erat, sed diam voluptua.

\begin{table}[H]\scriptsize
\centering
\begin{tabular}{c|cccc}
     Lines& F1 score&Validation Accuracy&MAP &Precision\\
    \hline
    Line\_1	& 91.2\,\%	& 63.7\,\%	& 0.72  &90.8\\
    Line\_2	& 78.3\,\%	& 63.7\,\%	& 0.83 &79.9\\
    Line\_3	& 77.1\,\%	& 61.2\,\%	& 0.73  &76.9\\
    Line\_4 & 89.9\,\%	& 63.7\,\%	& 0.74  &90.4\\
    Line\_5 & 80.7\,\% & 63.7\,\%  & 0.75 &81.6\\ 
\end{tabular}
\caption{Duis autem vel eum iriure dolor in hendrerit in vulputate velit esse molestie consequat, vel illum dolore eu feugiat nulla facilisis at vero eros et \acrfull{acronym3} accumsan et iusto odio dignissim qui blandit praesent luptatum zzril delenit augue duis dolore te feugait nulla facilisi.}
\label{tab:results}
\end{table}

Duis autem vel eum iriure dolor in hendrerit in vulputate velit esse molestie consequat, vel illum dolore eu \acrfull{acronym2} feugiat nulla facilisis at vero eros et accumsan et iusto odio dignissim qui blandit praesent luptatum zzril delenit augue duis dolore te feugait nulla facilisi. Lorem ipsum dolor sit amet, consectetuer adipiscing elit, sed diam nonummy nibh euismod tincidunt ut laoreet dolore magna aliquam erat volutpat. 

Ut wisi enim ad minim veniam, quis nostrud exerci tation ullamcorper suscipit lobortis nisl ut aliquip ex ea commodo consequat. Duis autem vel eum iriure dolor in hendrerit in vulputate velit esse molestie consequat, vel illum dolore eu feugiat nulla facilisis at vero eros et \acrfull{acronym3} accumsan et iusto odio dignissim qui blandit praesent luptatum zzril delenit augue duis dolore te feugait nulla facilisi.


\chapter{Chapter\_5 (optional)}

Lorem ipsum dolor sit amet \acrfull{acronym1}, consetetur sadipscing elitr, sed diam nonumy eirmod tempor invidunt ut labore et dolore magna aliquyam erat, sed diam voluptua. At vero eos et accusam et justo duo dolores et ea rebum. Stet clita kasd gubergren, no sea takimata sanctus est Lorem ipsum dolor sit amet. Lorem ipsum dolor sit amet, consetetur sadipscing elitr, sed diam nonumy eirmod tempor invidunt ut labore et dolore magna aliquyam erat, sed diam voluptua. At vero eos et accusam et justo duo dolores et ea rebum. Stet clita kasd gubergren, no sea takimata sanctus est Lorem ipsum dolor sit amet. Lorem ipsum dolor sit amet, consetetur sadipscing elitr, sed diam nonumy eirmod tempor invidunt ut labore et dolore magna aliquyam erat, sed diam voluptua. At vero eos et accusam et justo duo dolores et ea rebum. Stet clita kasd gubergren, no sea takimata sanctus est Lorem ipsum dolor sit amet. 

Duis autem vel eum iriure dolor in hendrerit in vulputate velit esse molestie consequat, vel illum dolore eu \acrfull{acronym2} feugiat nulla facilisis at vero eros et accumsan et iusto odio dignissim qui blandit praesent luptatum zzril delenit augue duis dolore te feugait nulla facilisi. Lorem ipsum dolor sit amet, consectetuer adipiscing elit, sed diam nonummy nibh euismod tincidunt ut laoreet dolore magna aliquam erat volutpat. 

Ut wisi enim ad minim veniam, quis nostrud exerci tation ullamcorper suscipit lobortis nisl ut aliquip ex ea commodo consequat. Duis autem vel eum iriure dolor in hendrerit in vulputate velit esse molestie consequat, vel illum dolore eu feugiat nulla facilisis at vero eros et \acrfull{acronym3} accumsan et iusto odio dignissim qui blandit praesent luptatum zzril delenit augue duis dolore te feugait nulla facilisi.

\section{Section\_1}

Lorem ipsum dolor sit amet \acrfull{acronym1}, consetetur sadipscing elitr, sed diam nonumy eirmod tempor invidunt ut labore et dolore magna aliquyam erat, sed diam voluptua. At vero eos et accusam et justo duo dolores et ea rebum. Stet clita kasd gubergren, no sea takimata sanctus est Lorem ipsum dolor sit amet. Lorem ipsum dolor sit amet, consetetur sadipscing elitr, sed diam nonumy eirmod tempor invidunt ut labore et dolore magna aliquyam erat, sed diam voluptua. At vero eos et accusam et justo duo dolores et ea rebum. Stet clita kasd gubergren, no sea takimata sanctus est Lorem ipsum dolor sit amet. Lorem ipsum dolor sit amet, consetetur sadipscing elitr, sed diam nonumy eirmod tempor invidunt ut labore et dolore magna aliquyam erat, sed diam voluptua.

\begin{table}[H]\scriptsize
\centering
\begin{tabular}{c|cccc}
     Lines& F1 score&Validation Accuracy&MAP &Precision\\
    \hline
    Line\_1	& 91.2\,\%	& 63.7\,\%	& 0.72  &90.8\\
    Line\_2	& 78.3\,\%	& 63.7\,\%	& 0.83 &79.9\\
    Line\_3	& 77.1\,\%	& 61.2\,\%	& 0.73  &76.9\\
    Line\_4 & 89.9\,\%	& 63.7\,\%	& 0.74  &90.4\\
    Line\_5 & 80.7\,\% & 63.7\,\%  & 0.75 &81.6\\ 
\end{tabular}
\caption{Duis autem vel eum iriure dolor in hendrerit in vulputate velit esse molestie consequat, vel illum dolore eu feugiat nulla facilisis at vero eros et \acrfull{acronym3} accumsan et iusto odio dignissim qui blandit praesent luptatum zzril delenit augue duis dolore te feugait nulla facilisi.}
\label{tab:eval}
\end{table}

Duis autem vel eum iriure dolor in hendrerit in vulputate velit esse molestie consequat, vel illum dolore eu \acrfull{acronym2} feugiat nulla facilisis at vero eros et accumsan et iusto odio dignissim qui blandit praesent luptatum zzril delenit augue duis dolore te feugait nulla facilisi. Lorem ipsum dolor sit amet, consectetuer adipiscing elit, sed diam nonummy nibh euismod tincidunt ut laoreet dolore magna aliquam erat volutpat. 

Ut wisi enim ad minim veniam, quis nostrud exerci tation ullamcorper suscipit lobortis nisl ut aliquip ex ea commodo consequat. Duis autem vel eum iriure dolor in hendrerit in vulputate velit esse molestie consequat, vel illum dolore eu feugiat nulla facilisis at vero eros et \acrfull{acronym3} accumsan et iusto odio dignissim qui blandit praesent luptatum zzril delenit augue duis dolore te feugait nulla facilisi.


\chapter{Summary and Outlook}

Lorem ipsum dolor sit amet \acrfull{acronym1}, consetetur sadipscing elitr, sed diam nonumy eirmod tempor invidunt ut labore et dolore magna aliquyam erat, sed diam voluptua. At vero eos et accusam et justo duo dolores et ea rebum. Stet clita kasd gubergren, no sea takimata sanctus est Lorem ipsum dolor sit amet. Lorem ipsum dolor sit amet, consetetur sadipscing elitr, sed diam nonumy eirmod tempor invidunt ut labore et dolore magna aliquyam erat, sed diam voluptua. At vero eos et accusam et justo duo dolores et ea rebum. Stet clita kasd gubergren, no sea takimata sanctus est Lorem ipsum dolor sit amet. Lorem ipsum dolor sit amet, consetetur sadipscing elitr, sed diam nonumy eirmod tempor invidunt ut labore et dolore magna aliquyam erat, sed diam voluptua. At vero eos et accusam et justo duo dolores et ea rebum. Stet clita kasd gubergren, no sea takimata sanctus est Lorem ipsum dolor sit amet. 

Duis autem vel eum iriure dolor in hendrerit in vulputate velit esse molestie consequat, vel illum dolore eu \acrfull{acronym2} feugiat nulla facilisis at vero eros et accumsan et iusto odio dignissim qui blandit praesent luptatum zzril delenit augue duis dolore te feugait nulla facilisi. Lorem ipsum dolor sit amet, consectetuer adipiscing elit, sed diam nonummy nibh euismod tincidunt ut laoreet dolore magna aliquam erat volutpat. 

Ut wisi enim ad minim veniam, quis nostrud exerci tation ullamcorper suscipit lobortis nisl ut aliquip ex ea commodo consequat. Duis autem vel eum iriure dolor in hendrerit in vulputate velit esse molestie consequat, vel illum dolore eu feugiat nulla facilisis at vero eros et \acrfull{acronym3} accumsan et iusto odio dignissim qui blandit praesent luptatum zzril delenit augue duis dolore te feugait nulla facilisi.

%======================================================================
%	Glossary
%======================================================================

%\manualmark
%\addcontentsline{table of Contents}{chapter}{Glossary}
%\markboth{Glossary}{Glossary}
%\def\glossaryname{Glossary}
%\printglossary
%\cleardoublepage

%======================================================================
%	Bibliography
%======================================================================
\markboth{Bibliography}{Bibliography}
\bibliographystyle{user}
\bibliography{literature}
%\cleardoublepage

%======================================================================
%	Theses
%======================================================================
%\chapter{Theses}
%\include{theses}


%======================================================================
%	Attachment
%======================================================================

\renewcommand\thechapter{\Alph{chapter}}
\setcounter{chapter}{0}
\renewcommand*{\theHchapter}{Appendix.\the\value{chapter}}
\chapter{Appendix\_A}
%\cleardoublepage
\label{chap:appendix}	

Lorem ipsum dolor sit amet \acrfull{acronym1}, consetetur sadipscing elitr, sed diam nonumy eirmod tempor invidunt ut labore et dolore magna aliquyam erat, sed diam voluptua. At vero eos et accusam et justo duo dolores et ea rebum. Stet clita kasd gubergren, no sea takimata sanctus est Lorem ipsum dolor sit amet. Lorem ipsum dolor sit amet, consetetur sadipscing elitr, sed diam nonumy eirmod tempor invidunt ut labore et dolore magna aliquyam erat, sed diam voluptua. At vero eos et accusam et justo duo dolores et ea rebum. Stet clita kasd gubergren, no sea takimata sanctus est Lorem ipsum dolor sit amet. Lorem ipsum dolor sit amet, consetetur sadipscing elitr, sed diam nonumy eirmod tempor invidunt ut labore et dolore magna aliquyam erat, sed diam voluptua. At vero eos et accusam et justo duo dolores et ea rebum. Stet clita kasd gubergren, no sea takimata sanctus est Lorem ipsum dolor sit amet. 

Duis autem vel eum iriure dolor in hendrerit in vulputate velit esse molestie consequat, vel illum dolore eu \acrfull{acronym2} feugiat nulla facilisis at vero eros et accumsan et iusto odio dignissim qui blandit praesent luptatum zzril delenit augue duis dolore te feugait nulla facilisi. Lorem ipsum dolor sit amet, consectetuer adipiscing elit, sed diam nonummy nibh euismod tincidunt ut laoreet dolore magna aliquam erat volutpat. 

Ut wisi enim ad minim veniam, quis nostrud exerci tation ullamcorper suscipit lobortis nisl ut aliquip ex ea commodo consequat. Duis autem vel eum iriure dolor in hendrerit in vulputate velit esse molestie consequat, vel illum dolore eu feugiat nulla facilisis at vero eros et \acrfull{acronym3} accumsan et iusto odio dignissim qui blandit praesent luptatum zzril delenit augue duis dolore te feugait nulla facilisi.

%======================================================================
%	Declaration of Authorship
%======================================================================
\chapter*{Declaration of Authorship}
\thispagestyle{empty}

I hereby certify that this research project has been composed by me and is based on my own work, unless stated otherwise. No other person's work has been used without due acknowledgement in this project. All references and verbatim extracts have been quoted, and all sources of information, including graphs and data sets, have been specifically acknowledged. \\[2ex]
\dcplace, \dcdate\\[6ex]
\flushleft
\newlength\us
\settowidth{\us}{-\dcauthorfirstname~\dcauthorlastname-}
\begin{tabular}{p{\us}}\hline
\centering\footnotesize \dcauthorfirstname~\dcauthorlastname
\end{tabular}



\end{document}
